% \let must keep the engine's macro table honest in both directions: binding a
% name that *was* a macro to something that is not one has to drop the old
% body, or the stale entry wins on every later lookup.
\input plain
\catcode`\@=11
\let\bar= X
\message{<char=\meaning\bar>}
\def\m{OLD}\let\m\relax        \message{<prim=\meaning\m>}
\def\n{OLD}\let\n\undefinedcs  \message{<undef=\meaning\n>}
\def\p{OLD}\def\q{NEW}\let\p\q \message{<macro=\meaning\p>}
\def\r{OLD}\let\r\count        \message{<register=\meaning\r>}
\def\s{OLD}\chardef\ch=65 \let\s\ch \message{<chardef=\meaning\s>}
% Binding a space token: \let skips one optional space after the `=`, so a
% second one is what gets bound (tex.web 1221). expl3 defines \exp_stop_f:
% exactly this way, via `\let #1 =~` with the argument supplying the second.
\def\setsp#1{\let#1= }
\setsp\spacetok\relax
\message{<space=\meaning\spacetok>}
% And the round trip: a name let to a macro is still a macro.
\def\t{TEE}\let\u\t \message{<roundtrip=\u|\meaning\u>}
\bye
